\documentclass[12pt,a4paper]{article}
\usepackage[UTF8]{ctex}
\usepackage{amsmath, amssymb}
\usepackage{geometry}
\usepackage{booktabs}
\usepackage{graphicx}
\usepackage{float}
\usepackage{multirow}
\usepackage{enumitem}
\usepackage{hyperref}

\geometry{left=2.5cm, right=2.5cm, top=2.5cm, bottom=2.5cm}

\title{\textbf{数据分析（论文第 X 章素材）}\\[6pt]
\large 数据总览 + 四个问题分别的数据分析}
\author{2026 年第七届"华数杯"大学生数学建模竞赛 B 题}
\date{2026/08/10}

\begin{document}
\maketitle

% ============================================================
\section{数据总览与建模动机}
% ============================================================

本节先给出三组芯片的总体数据画像，再从\textbf{模块几何特征、网表连接
结构、搜索空间规模}三个维度给出与建模直接相关的结论；后续四节分别
围绕各问题的数据展开分析。

\subsection{三芯片规模总览}

\begin{table}[H]
\centering
\scriptsize
\setlength{\tabcolsep}{3pt}
\caption{三组芯片原始数据规模}
\begin{tabular}{lccccc}
\toprule
芯片 & 模块数 & 端子数 & 线网数 & 模块总面积 & 面积利用率(基准轮廓) \\
\midrule
n100 & 100 & 334 & 885 & 179501 & 0.9487（自由轮廓）/ 0.867（$d{=}0.15$） \\
n200 & 200 & 564 & 1585 & 175696 & 0.9595 / 0.8676 \\
n300 & 300 & 569 & 1893 & 273170 & 0.9577 / 0.868 \\
\bottomrule
\end{tabular}
\end{table}

\subsection{三维度总体结论}

\begin{enumerate}[leftmargin=2em]
    \item \textbf{几何异质性}（服务于问题一、四）：面积变异系数
    $CV \approx 0.5$、面积跨度 11.6--14.5 倍、极端长条模块
    （长宽比 $>3$）占 3.3\%--6.0\%——中等偏上的异质性意味着
    随机摆放易在局部留下不可填补死区；
    \item \textbf{网表连接集中性}（服务于问题二）：节点度均值
    12--15、最大 28--35，存在明确的枢纽模块（Hub），但双对数坐标下
    分布\textbf{非严格幂律}（拟合 $R^2 \le 0.63$）——集中型连接而非
    无标度网络；
    \item \textbf{搜索空间规模}（服务于问题一、二、三的算法选择）：
    $|S| = n!\cdot C_n \cdot 2^n$，$n=300$ 时 $\log_{10}|S| \approx
    881$（约 $10^{881}$，远超宇宙原子数 $10^{80}$）——精确算法不可行，
    必须采用快速模拟退火。
\end{enumerate}

% ============================================================
\section{问题一的数据分析：自由轮廓面积最小化}
% ============================================================

\subsection{模块几何异质性}

\begin{table}[H]
\centering
\scriptsize
\setlength{\tabcolsep}{3pt}
\caption{问题一相关几何指标}
\begin{tabular}{lcccc}
\toprule
芯片 & 面积 CV & 面积跨度 $\max/\min$ & 长宽比 $>3$ 数量 & 长宽比 $>3$ 占比 \\
\midrule
n100 & 0.498 & 11.6$\times$ & 6 & 6.0\% \\
n200 & 0.542 & 14.5$\times$ & 8 & 4.0\% \\
n300 & 0.528 & 14.5$\times$ & 10 & 3.3\% \\
\bottomrule
\end{tabular}
\end{table}

\begin{figure}[H]
\centering
\includegraphics[width=0.95\textwidth]{../图表/公共/数据总览/eda_09_area_kde.png}
\caption{模块面积核密度估计（$\log_{10}$ 横轴）：右偏长尾分布}
\label{fig:q1_kde}
\end{figure}

\begin{figure}[H]
\centering
\includegraphics[width=0.95\textwidth]{../图表/公共/数据总览/eda_10_wvsh_scatter.png}
\caption{模块宽高散点（灰虚线 $W=H$，红圈为长宽比 $>3$ 的极端长条模块）}
\label{fig:q1_scatter}
\end{figure}

\subsection{结论话术（建议直接采用）}

"问题一要求在自由轮廓下最小化包围盒面积。数据表明三芯片模块面积
分布的变异系数约 0.5，面积跨度达 11.6--14.5 倍，并存在 3.3\%--6.0\%
的极端长条形模块。这种\textbf{中等偏上的几何异质性}意味着：若随机
摆放，大模块与长条模块周围极易形成小模块无法填补的局部空隙，包围盒
面积将显著偏离模块总面积下界。这从数据上解释了为何需要\textbf{基于
B$^*$ 树编码}的紧凑表示（解码天然压缩）与\textbf{多扰动算子
（旋转/删除-插入/交换）+ 多起点}的搜索策略：旋转算子专门应对异质
长宽比，删除-插入负责大范围重排，多起点对冲局部最优。"

% ============================================================
\section{问题二的数据分析：固定轮廓 HPWL 最小化}
% ============================================================

\subsection{网表连接结构}

\begin{table}[H]
\centering
\scriptsize
\setlength{\tabcolsep}{3pt}
\caption{问题二相关网表指标}
\begin{tabular}{lccccc}
\toprule
芯片 & 线网数 & 模块数 & 度均值 & 度最大值 & 度范围 \\
\midrule
n100 & 885 & 100 & 15.4 & 28 & 6--28 \\
n200 & 1585 & 200 & 15.2 & 35 & 6--35 \\
n300 & 1893 & 300 & 12.6 & 35 & 6--35 \\
\bottomrule
\end{tabular}
\end{table}

\begin{table}[H]
\centering
\scriptsize
\setlength{\tabcolsep}{3pt}
\caption{连接度最高的 5\% 枢纽模块}
\begin{tabular}{llllll}
\toprule
芯片 & \multicolumn{5}{c}{枢纽模块（模块:度）} \\
\midrule
n100 & b73:28 & b66:28 & b49:27 & b44:27 & b15:27 \\
n200 & b21:35 & b28:35 & b64:29 & b143:28 & b129:28 \\
n300 & b57:35 & b33:24 & b131:23 & b266:23 & b92:23 \\
\bottomrule
\end{tabular}
\end{table}

\begin{figure}[H]
\centering
\includegraphics[width=0.95\textwidth]{../图表/公共/数据总览/eda_11_degree_loglog.png}
\caption{节点度数分布（双对数坐标）：集中型，非幂律}
\label{fig:q2_deg}
\end{figure}

\begin{figure}[H]
\centering
\includegraphics[width=0.95\textwidth]{../图表/公共/数据总览/eda_12_hubs.png}
\caption{Top 5\% 枢纽模块连接度}
\label{fig:q2_hub}
\end{figure}

\subsection{结论话术（建议直接采用）}

"问题二在固定轮廓内最小化总半周长线长（HPWL）。将网表抽象为超图后
统计显示：节点度均值仅 12--15，而最高连接模块度达 28--35，存在明确
的枢纽模块。双对数坐标下分布\textbf{非严格幂律}（$R^2 \le 0.63$），
属\textbf{集中型连接}——少数高连模块承载了不成比例的网络连接。在
HPWL 目标下，将高连模块布局于相互邻近区域可显著压缩其参与线网的
引脚跨度，因此\textbf{对高连模块的拓扑位置施加引导}具有明确收益
（本文当前版本未引入该引导，列为后续改进方向）。此外轮廓死区
$d=0.15$ 使面积利用率存在上界 $1/1.15=0.8696$，实测 0.867--0.868
已逼近该极限，说明布局紧凑度不是主要矛盾，\textbf{线长优化}才是
问题二的核心。"

% ============================================================
\section{问题三的数据分析：最小可行死区比例 $d^*$}
% ============================================================

\subsection{轮廓与利用率下界}

\begin{table}[H]
\centering
\scriptsize
\setlength{\tabcolsep}{3pt}
\caption{问题三相关指标（轮廓 $\text{side}(d)=\lceil\sqrt{\Sigma A(1+d)}\rceil$）}
\begin{tabular}{lcccccc}
\toprule
芯片 & 模块总面积 $T$ & side($d{=}0$) & side($d{=}0.15$) & 利用率上界 $1/1.15$ & 基准利用率 & 定稿 $d^*$ \\
\midrule
n100 & 179501 & 424 & 455 & 0.8696 & 0.867 & 0.076072 \\
n200 & 175696 & 420 & 450 & 0.8696 & 0.8676 & 0.084141 \\
n300 & 273170 & 523 & 561 & 0.8696 & 0.868 & 0.111773 \\
\bottomrule
\end{tabular}
\end{table}

\begin{figure}[H]
\centering
\includegraphics[width=0.95\textwidth]{../图表/公共/数据总览/eda_06_packing.png}
\caption{面积利用率预分析（自由轮廓下模块总面积的铺满潜力）}
\label{fig:q3_pack}
\end{figure}

\subsection{结论话术（建议直接采用）}

"问题三要求寻找最小的可行死区比例 $d^*$。数据表明：三芯片模块总
面积 $T$ 确定后，轮廓边长 $\text{side}(d)$ 随 $d$ 单调不减，$d=0$
时轮廓即面积下界（n300 为 $523\times523$），$d=0.15$ 时利用率上界
为 $1/1.15=0.8696$。自由轮廓下的布局利用率实测高达 0.95 左右
（问题一结果），说明\textbf{模块总面积的铺满潜力可观}；而固定轮廓
下的死区比例从 $d=0.15$ 向下压缩的空间，正是问题三要探测的临界
区域。数据支持将该问题形式化为\textbf{单调可行性判定 + 二分搜索}，
并通过多种子判定抑制启发式假阴性。"

% ============================================================
\section{问题四的数据分析：L/T 模块最小包围盒}
% ============================================================

\subsection{模块几何：矩形分解与凹角}

\begin{table}[H]
\centering
\scriptsize
\setlength{\tabcolsep}{3pt}
\caption{问题四模块几何（矩形分解表示）}
\begin{tabular}{lcllll}
\toprule
模块 & 形状 & 分解矩形（相对坐标） & 面积 & bbox 面积 & 凹角面积 \\
\midrule
b1 & T 型 & 顶横条 $4\times2$@(0,2)；竖干 $2\times2$@(1,0) & 12 & 16 & 4（25\%） \\
b2 & L 型 & 左竖条 $1\times4$@(0,0)；右下横条 $1\times2$@(1,0) & 6 & 8 & 2（25\%） \\
b3 & 矩形 & $2\times1$ & 2 & 2 & 0 \\
b4 & 矩形 & $1\times4$ & 4 & 4 & 0 \\
\midrule
合计 & & & 24 & 30 & 6 \\
\bottomrule
\end{tabular}
\end{table}

\subsection{结论话术（建议直接采用）}

"问题四引入 L/T 型异形模块并允许四向旋转。数据表明：T 型模块 b1
（面积 12）的包围盒为 $4\times4$（面积 16），L 型模块 b2（面积 6）
的包围盒为 $2\times4$（面积 8）——\textbf{凹角空缺分别占包围盒的
25\%}。若按包围盒整体放置（bbox 版本），6 个模块的包围盒面积合计
30，而实际总面积仅 24，凹角空缺将直接转化为死区。这从数据上支持
\textbf{矩形分解 + 凹角精确支撑}的建模方式：将模块表示为固定相对
位置的矩形并集，放置时以内部矩形逐一对轮廓求支撑，可复用凹角空间；
模块总面积 24 构成包围盒面积的\textbf{严格下界}，穷举全部
86\,016 种组合并达到下界即可\textbf{证明全局最优}。"

% ============================================================
\section{小结}

四个问题的数据分别支撑各自的建模动机：问题一依赖几何异质性分析
（多扰动 + 多起点）、问题二依赖网表集中性分析（HPWL 主导 + 高连
模块引导展望）、问题三依赖轮廓与利用率下界分析（单调判定 + 二分）、
问题四依赖凹角空缺量化分析（矩形分解 + 下界可达最优性证明）。

\end{document}
